\documentclass[11pt,a4paper]{article}
\usepackage[T1]{fontenc}
\usepackage[utf8]{inputenc}
\usepackage{lmodern}
\usepackage{amsmath,amssymb,amsthm}
\usepackage{enumitem}
\usepackage{geometry}
\geometry{margin=2.5cm}
\setlist[itemize]{noitemsep,topsep=2pt}
\newcommand{\E}{\mathbb{E}}
\newcommand{\Pp}{\mathbb{P}}
\title{Proof summaries}
\author{}
\date{}
\begin{document}
\maketitle
\noindent\textbf{How to use.} For each theorem, the bullet points are the ``story'' of the proof: (i) main tool/lemma, (ii) key decomposition/induction, (iii) conclusion/identification.\par\bigskip
\subsection*{Theorem 1.6 (Conditional expectation existence)}
\textbf{Proof idea.}
\begin{itemize}
\item Define a finite measure $Q$ on $(\Omega,\mathcal F)$ by $Q(F)=\int_F X\,d\mathbb P$ (first for $X\ge0$).
\item Note $Q\ll \mathbb P|_{\mathcal F}$, so Radon--Nikodym gives an $\mathcal F$-measurable density $Z=\frac{dQ}{d\mathbb P|_{\mathcal F}}$ with $\int_F Z\,d\mathbb P=\int_F X\,d\mathbb P$ for all $F\in\mathcal F$.
\item For general $X\in L^1$, apply the nonnegative case to $X^+$ and $X^-$ and set $Z=Z^+-Z^-$.
\item Uniqueness: if $Z,Z'$ both work, test on $\{Z-Z'>1/n\}\in\mathcal F$ to force $\mathbb P(Z>Z')=\mathbb P(Z'<Z)=0$.
\end{itemize}

\subsection*{Theorem 1.8 (Properties of conditional expectation)}
\textbf{Proof idea.}
\begin{itemize}
\item Fix versions $X_0=\mathbb E[X\mid\mathcal F]$, $Y_0=\mathbb E[Y\mid\mathcal F]$ and always prove identities by checking the defining property on every $F\in\mathcal F$ and then invoking a.s.\ uniqueness.
\item Linearity and constants: build the candidate (e.g. $aX_0+bY_0$ or the constant $c$), verify it is $\mathcal F$-measurable, and match integrals over $F\in\mathcal F$.
\item Monotonicity: consider $A=\{X_0>Y_0\}\in\mathcal F$ and integrate $Y_0-X_0$ over $A$; the sign contradiction forces $\mathbb P(A)=0$.
\item Tower property: test on $G\in\mathcal G\subset\mathcal F$ and use the defining identity twice.
\item Measurable factor / independence: prove first for indicators $Z=\mathbf 1_{F_0}$, extend to simple functions and then to general $Z$ by dominated convergence; for independence test $\int_F X\,d\mathbb P=\mathbb E[X]\mathbb P(F)$.
\end{itemize}

\subsection*{Theorem 1.9 (Best approximation / orthogonal projection)}
\textbf{Proof idea.}
\begin{itemize}
\item Let $X_0=\mathbb E[X\mid\mathcal F]$ and take any $Y\in L^2(\Omega,\mathcal F,\mathbb P)$.
\item Show orthogonality: $\mathbb E[(X-X_0)(X_0-Y)]=\mathbb E\big[\mathbb E[X-X_0\mid\mathcal F](X_0-Y)\big]=0$ (tower + $\mathcal F$-measurable factor).
\item Apply Pythagoras in $L^2$: $\|X-Y\|_2^2=\|X-X_0\|_2^2+\|X_0-Y\|_2^2\ge \|X-X_0\|_2^2$.
\item Equality iff $\|X_0-Y\|_2=0$, i.e. $Y=X_0$ a.s., so $X_0$ is the orthogonal projection / best $L^2$ predictor.
\end{itemize}

\subsection*{Theorem 1.10 (Jensen inequality)}
\textbf{Proof idea.}
\begin{itemize}
\item Use the supporting-line characterization of convexity: for each $x_0$, $g(x)\ge g(x_0)+m_{x_0}(x-x_0)$ for some slope $m_{x_0}$ (e.g. right derivative).
\item Choose a countable set of supporting lines (take $x_0\in I\cap\mathbb Q$) so that $g$ is the pointwise supremum of these affine functions.
\item Insert $X$ and take $\mathbb E[\cdot\mid\mathcal F]$: conditional expectation preserves affine functions, so each supporting-line inequality becomes an inequality involving $\mathbb E[X\mid\mathcal F]$.
\item Take the supremum over the countable family (so no measurability issues beyond a single null set) to get $g(\mathbb E[X\mid\mathcal F])\le \mathbb E[g(X)\mid\mathcal F]$ a.s.
\end{itemize}

\subsection*{Theorem 2.4 (Existence of an RCD for real-valued random variables)}
\textbf{Proof idea.}
\begin{itemize}
\item For each rational $r$, pick a version $F(r,\omega)=\mathbb P(Y\le r\mid\mathcal F)(\omega)$; these are $\mathcal F$-measurable numbers in $[0,1]$.
\item Use monotonicity of conditional expectation to enforce $r\mapsto F(r,\omega)$ nondecreasing outside one null set (take a countable union over $r\le s$ in $\mathbb Q$).
\item Force right limits and boundary conditions ($F(-n)\downarrow0$, $F(n)\uparrow1$) outside another null set using dominated convergence.
\item Define for all real $z$ the right-continuous hull $\tilde F(z,\omega)=\inf\{F(r,\omega):r\in\mathbb Q, r>z\}$; for each good $\omega$ this is a genuine cdf.
\item Let $\vartheta(\omega,\cdot)$ be the probability measure with cdf $\tilde F(\cdot,\omega)$; measurability on the generating class $(-\infty,r]$ and a $\pi$--$\lambda$ (monotone class) argument give measurability for all Borel sets.
\item Verify the defining identity first on $(-\infty,r]$ and extend to all $B\in\mathcal B(\mathbb R)$ by uniqueness of measures on generators.
\end{itemize}

\subsection*{Theorem 2.8 (Polish $\implies$ Borel space)}
\textbf{Proof idea.}
\begin{itemize}
\item Use that a Polish space has a countable base (separability) and nice limit behavior (completeness), which yields a standard Borel structure.
\item Concretely, one can embed $E$ as a Borel subset of a universal Polish space (e.g. $[0,1]^{\mathbb N}$) via a Borel isomorphism built from a countable base / separating family.
\item Thus $(E,\mathcal B(E))$ is measurably isomorphic to a Borel subset of $\mathbb R$, i.e. a Borel (standard Borel) space.
\end{itemize}

\subsection*{Theorem 2.10 (Existence of RCDs for Borel-valued random variables)}
\textbf{Proof idea.}
\begin{itemize}
\item Take a measurable isomorphism $\xi:(E,\mathcal E)\to (B,\mathcal B(B))$ with $B\subset\mathbb R$ Borel (exists since $E$ is a Borel space).
\item Push $Y$ forward to a real-valued variable $\tilde Y=\xi\circ Y$ and apply the real-valued existence theorem to get an RCD $\vartheta_{\tilde Y,\mathcal F}$.
\item Pull the kernel back: define $\vartheta_{Y,\mathcal F}(\omega,A)=\vartheta_{\tilde Y,\mathcal F}(\omega,\xi(A))$; measurability and the defining identity transfer through the isomorphism.
\end{itemize}

\subsection*{Theorem 2.11 (Integral representation of conditional expectation)}
\textbf{Proof idea.}
\begin{itemize}
\item First prove for simple $f=\sum_i \alpha_i\mathbf 1_{A_i}$ using the RCD property: $\mathbb E[f(X)\mathbf 1_B]=\sum_i \alpha_i\mathbb P(B\cap\{X\in A_i\})=\int_B\int f\,d\vartheta_{X,\mathcal F}\,d\mathbb P$.
\item Approximate any nonnegative $f$ by an increasing sequence of simple functions and pass to the limit via monotone convergence (both inside the kernel integral and in expectation).
\item Recognize the resulting $\mathcal F$-measurable function $\omega\mapsto\int f\,d\vartheta_{X,\mathcal F}(\omega,\cdot)$ as the unique version of $\mathbb E[f(X)\mid\mathcal F]$ by testing against indicators of $B\in\mathcal F$.
\end{itemize}

\subsection*{Theorem 2.12 (Disintegration lemma for probability measures on products)}
\textbf{Proof idea.}
\begin{itemize}
\item Work on the probability space $(E\times F,\mathcal E\otimes\mathcal F,\mu)$ with coordinate projections $X(x,y)=x$ and $Y(x,y)=y$.
\item By existence of RCDs on Borel spaces, take an RCD of $Y$ given $\sigma(X)$; it is a kernel $\vartheta_{Y,\sigma(X)}(\omega,\cdot)$.
\item Factor the $\sigma(X)$-measurable map $\omega\mapsto\vartheta_{Y,\sigma(X)}(\omega,B)$ through $X$ to get a measurable $\kappa(x,B)$ (use the factorization lemma).
\item Then $\kappa(x,\cdot)$ is a probability measure for $\mu_E$-a.e.\ $x$; redefine on a null set to get it for all $x$.
\item Finally, the RCD identity on sets $A\times F$ yields $\mu(A\times B)=\int_A \kappa(x,B)\,\mu_E(dx)$, and uniqueness follows from agreement on a countable generator + monotone class.
\end{itemize}

\subsection*{Theorem 3.7 (Snell envelope recursion and existence of an optimal rule)}
\textbf{Proof idea.}
\begin{itemize}
\item Backward induction on time: at $N$ you must stop, so $U_N=Z_N$.
\item For the recursion, fix $T\in\mathcal T_{n-1}$ and split on $\{T=n-1\}$ vs. $\{T\ge n\}$; on the second event use the tower property and the definition of $U_n$ to bound $\mathbb E[Z_T\mid\mathcal F_{n-1}]$ by $\max\{Z_{n-1},\mathbb E[U_n\mid\mathcal F_{n-1}]\}$.
\item Taking the essential supremum over $T$ gives $U_{n-1}\le$ RHS; equality comes from considering the two candidate strategies: stop now or continue optimally from $n$.
\item Define $T_n^*=\min\{k\ge n:U_k=Z_k\}$; measurability of $\{T_n^*=k\}$ follows from adaptedness of $U,Z$.
\item Optimality is proved by induction using the recursion: on the continuation set $\{T_n^*\ge n+1\}$ you have $U_n=\mathbb E[U_{n+1}\mid\mathcal F_n]$, and eventually $U_n=\mathbb E[Z_{T_n^*}\mid\mathcal F_n]$.
\end{itemize}

\subsection*{Theorem 3.10 (Dynamic programming / Bellman recursion)}
\textbf{Proof idea.}
\begin{itemize}
\item Use induction to show the Snell envelope has the Markov form $U_n=V_n(X_n)$.
\item Assume $U_n=V_n(X_n)$; then by the Snell recursion $U_{n-1}=\max\{g_{n-1}(X_{n-1}),\mathbb E[V_n(X_n)\mid\mathcal F_{n-1}]\}$.
\item By the Markov property, $\mathbb E[V_n(X_n)\mid\mathcal F_{n-1}]=\mathbb E[V_n(X_n)\mid X_{n-1}]=\int_E V_n(y)Q(X_{n-1},dy)$, which defines $V_{n-1}$.
\item The optimal stopping time is the first entry into the stopping region $\{x:g_n(x)=V_n(x)\}$ (same hitting-time idea as in Theorem 3.7).
\end{itemize}

\subsection*{Theorem 3.14 (Martingale optimality principle)}
\textbf{Proof idea.}
\begin{itemize}
\item Part (i) is the minimality property of the Snell envelope: by definition it dominates $Z$ and is the smallest supermartingale doing so.
\item For any stopping time $T$, optional sampling for supermartingales gives $\mathbb E[Z_T]\le \mathbb E[U_T]\le \mathbb E[U_0]$, proving (ii).
\item For $T^*=\min\{n:U_n=Z_n\}$, the Snell recursion implies that $U$ is a martingale up to $T^*$, hence $\mathbb E[U_{T^*}]=\mathbb E[U_0]$ and $U_{T^*}=Z_{T^*}$; therefore (iii).
\end{itemize}

\subsection*{Theorem 3.15 (Doob decomposition)}
\textbf{Proof idea.}
\begin{itemize}
\item Define the compensator increments $\Delta A_n := Y_{n-1}-\mathbb E[Y_n\mid\mathcal F_{n-1}]\ge0$; this is exactly the supermartingale drift.
\item Let $A_0=0$ and $A_n=\sum_{k=1}^n \Delta A_k$; then $A$ is predictable and increasing.
\item Set $M_n:=Y_n+A_n$ and check the martingale property by direct conditioning: $\mathbb E[M_n\mid\mathcal F_{n-1}]=Y_{n-1}+A_{n-1}=M_{n-1}$.
\item Thus $Y=M-A$; uniqueness follows from uniqueness of conditional expectations / increments.
\end{itemize}

\subsection*{Theorem 3.21 (Threshold rule and asymptotics)}
\textbf{Proof idea.}
\begin{itemize}
\item Reduce the secretary problem to a 1D Markov optimal stopping problem using the relative rank process $X_n$ and the Snell/Bellman recursion for success probabilities.
\item Show the continuation value $c_n=\mathbb E[U_{n+1}\mid\mathcal F^X_n]$ becomes an explicit average because $X_{n+1}$ is conditionally uniform on $\{1,\dots,n+1\}$ given past relative ranks (symmetry of random permutations).
\item The stopping rule compares `stop now' (only beneficial when $X_n=1$) versus `continue' (value $c_n$), yielding a threshold index $m$ where the inequality flips.
\item Asymptotics come from approximating sums by integrals, giving $m\approx N/e$ and success probability $\approx 1/e$.
\end{itemize}

\subsection*{Theorem 4.3 (Conditional Gaussianity: informal statement)}
\textbf{Proof idea.}
\begin{itemize}
\item Exploit linear-Gaussian dynamics: $(X_{n+1},Y_{n+1})$ is an affine function of $(X_n)$ plus independent Gaussian noise.
\item By induction, conditional on past observations $\mathcal Y_n$, the state $X_n$ remains Gaussian (Gaussian family closed under affine maps and conditioning).
\item Compute how the mean and covariance propagate: prediction via the linear dynamics, update via Gaussian conditioning (a regression formula).
\end{itemize}

\subsection*{Theorem 4.5 (Normal correlation / Gaussian regression with pseudo-inverse)}
\textbf{Proof idea.}
\begin{itemize}
\item Write the joint vector $(X,Y)$ as Gaussian with block covariance; compute the conditional mean as the $L^2$-best linear predictor (orthogonal projection in a Gaussian Hilbert space).
\item Use the normal equations to solve for the regression matrix; when the covariance is not invertible use the Moore--Penrose pseudo-inverse to get the minimal-norm solution.
\item Show the conditional covariance is the Schur complement (or its pseudo-inverse analogue), and verify the result by checking characteristic functions / completing the square in the Gaussian density.
\end{itemize}

\subsection*{Theorem 2.2 (Conditional Gaussianity)}
\textbf{Proof idea.}
\begin{itemize}
\item Set up the linear state-space model and show $(X_n,Y_n)$ is jointly Gaussian for each fixed $n$ by affine transformations of Gaussians.
\item Use that conditioning a multivariate Gaussian remains Gaussian with explicitly updated mean/covariance (given by regression).
\item Conclude that the filtering distribution $\mathcal L(X_n\mid\mathcal Y_n)$ is Gaussian and therefore fully described by $(m_n,\Sigma_n)$.
\end{itemize}

\subsection*{Theorem 5.4 (Equivalent disturbance representation)}
\textbf{Proof idea.}
\begin{itemize}
\item Unroll the controlled dynamics to express $X_n$ as a measurable function of initial state, policy decisions, and an i.i.d.\ disturbance sequence.
\item Construct the disturbance space as a product space and define the controlled process as a coordinate map; verify it has the same transition kernel as the original model.
\item This gives an equivalent `canonical' representation where randomness is separated (noise) from control (policy), which simplifies measurability and dynamic programming arguments.
\end{itemize}

\subsection*{Theorem 5.13 (History dependence does not improve the value)}
\textbf{Proof idea.}
\begin{itemize}
\item Compare a history-dependent policy with its `Markovized' version: at each time $n$ pick an action distribution depending only on the current state, equal to the conditional law of the original action given the state.
\item Show by induction that both policies generate the same state distribution at all times (use the tower property and the transition kernel).
\item Therefore they have the same expected total reward, so optimizing over Markov policies is sufficient.
\end{itemize}

\subsection*{Theorem 6.4 (Reward Iteration)}
\textbf{Proof idea.}
\begin{itemize}
\item Define reward iteration under a fixed policy by repeatedly applying the one-step operator $T_n^{f_n}$ backward from time $N$.
\item Prove by induction on $n$ that the iterates equal the conditional expected remaining reward starting from state $x$ at time $n$.
\item The key step is the tower property + Markov property to replace conditioning on full history by conditioning on the current state.
\end{itemize}

\subsection*{Theorem 6.6 (Verification Theorem)}
\textbf{Proof idea.}
\begin{itemize}
\item Take a candidate function sequence $(v_n)$ satisfying the Bellman inequalities $v_n\ge T_nv_{n+1}$ and terminal condition; show by backward induction that $v_n$ dominates the value function $V_n$ (monotonicity of $T_n$).
\item If equality holds and maximizers exist, build a policy choosing stage-wise maximizers; then reward iteration shows this policy attains $v_n$.
\item Conclude $v_n=V_n$ and the constructed policy is optimal.
\end{itemize}

\subsection*{Theorem 6.8 (Structure Theorem)}
\textbf{Proof idea.}
\begin{itemize}
\item Backward induction using the structure assumption: start with $V_N=g_N\in M_N$.
\item Assume $V_{n+1}\in M_{n+1}$; then closure (SAN-ii) gives measurability of $T_nV_{n+1}$ and keeps you in $M_n$.
\item Upper bound: any policy value satisfies $V_n^\pi\le T_nV_{n+1}$, hence $V_n\le T_nV_{n+1}$ after taking the supremum.
\item Lower bound: pick a maximizer $f_n^*$ for $V_{n+1}$ (SAN-iii) and patch with $\varepsilon$-optimal continuation to show $V_n\ge T_nV_{n+1}$.
\item Collect the maximizers over all times and apply the verification theorem to get an optimal policy.
\end{itemize}

\subsection*{Theorem 7.15 (Daniell–Kolmogorov)}
\textbf{Proof idea.}
\begin{itemize}
\item Define a pre-measure on cylinder sets using the consistent finite-dimensional distributions.
\item Show this pre-measure is countably additive on the algebra of cylinder sets (consistency is crucial).
\item Apply Carathéodory’s extension theorem to extend uniquely to the product $\sigma$-algebra, obtaining a probability measure on the path space.
\item Finally, the coordinate maps on the product space form the desired process with the given finite-dimensional laws.
\end{itemize}

\subsection*{Theorem 7.18 (Existence of stochastic processes in Lusin spaces)}
\textbf{Proof idea.}
\begin{itemize}
\item Reduce the problem to the standard Borel/Polish case by representing the Lusin space as a Borel subset of a Polish space via a measurable isomorphism.
\item Apply Daniell--Kolmogorov on the Polish model space to get a process there.
\item Push the process forward through the isomorphism to obtain a process with the desired finite-dimensional distributions on the original Lusin space.
\end{itemize}

\subsection*{Theorem 8.1 (Existence (L´evy–Ciesielski / Haar–Schauder Construction)}
\textbf{Proof idea.}
\begin{itemize}
\item Choose an orthonormal basis of $L^2[0,1]$ with good path properties (Haar functions) and define $B_t$ as a random series with i.i.d.\ Gaussian coefficients (L\'evy--Ciesielski).
\item Use Kolmogorov’s three-series / $L^2$ estimates to show the series converges uniformly a.s., giving continuous sample paths.
\item Compute covariances to verify $\mathbb E[B_t]=0$ and $\mathrm{Cov}(B_s,B_t)=\min\{s,t\}$, hence the finite-dimensional distributions are those of Brownian motion.
\item Independence of increments follows from orthogonality of basis pieces corresponding to disjoint intervals.
\end{itemize}

\subsection*{Theorem 9.2 (Optional Stopping Theorem)}
\textbf{Proof idea.}
\begin{itemize}
\item First prove the inequality for bounded stopping times by writing $\mathbb E[M_T]=\sum_n \mathbb E[M_n\mathbf 1_{\{T=n\}}]$ and using the martingale property stepwise.
\item For general integrable/bounded-below martingales, approximate $T$ by $T\wedge n$ and pass to the limit using dominated/monotone convergence plus uniform integrability assumptions.
\item The key idea is: martingales have zero conditional drift, so stopping cannot create expected gain unless integrability fails.
\end{itemize}

\subsection*{Theorem 9.3 (Skorokhod embedding)}
\textbf{Proof idea.}
\begin{itemize}
\item Construct a stopping time $\tau$ as the first exit time of Brownian motion from a carefully chosen interval depending on the target law $\mu$.
\item Use the martingale $B_t$ and (optional stopping) on bounded approximations to ensure $\mathbb E[B_\tau]=0$ and to match second moments.
\item Choose the interval/randomization so that the distribution of $B_\tau$ matches $\mu$ exactly (typically via quantile or potential-theoretic arguments).
\item Uniform integrability / bounded stopping-time approximations justify passing limits and conclude $B_\tau\sim\mu$.
\end{itemize}

\subsection*{Theorem 9.12 (Strong Markov property)}
\textbf{Proof idea.}
\begin{itemize}
\item Use the exponential martingale $M_t(\theta)=\exp(i\theta B_t-\tfrac12\theta^2 t)$ and optional stopping at $T\wedge n$ to compute conditional characteristic functions of increments after $T$.
\item Let $n\to\infty$ and use continuity of paths + dominated convergence to pass from $T\wedge n$ to $T$.
\item The conditional characteristic function of $B_{T+t}-B_T$ given $\mathcal F_T$ equals $\exp(-\tfrac12\theta^2 t)$, hence the increment is $N(0,t)$.
\item Because this conditional characteristic function is deterministic, the increment is independent of $\mathcal F_T$; iterate to obtain independent stationary increments after $T$.
\end{itemize}

\subsection*{Theorem 9.14 (Reflection principle)}
\textbf{Proof idea.}
\begin{itemize}
\item For $a>0$, relate paths that hit level $a$ before time $t$ to paths whose endpoint exceeds $a$ by reflecting the trajectory after the first hitting time $\tau_a$.
\item Show the reflection map preserves Wiener measure and swaps the events $\{\sup_{s\le t}B_s\ge a,\ B_t<b\}$ and $\{B_t>2a-b\}$.
\item This yields $\mathbb P(\sup_{s\le t}B_s\ge a)=2\mathbb P(B_t\ge a)$ and more general joint identities for $(\sup B_s, B_t)$.
\end{itemize}

\subsection*{Theorem 10.18 (Relation E to M (E)}
\textbf{Proof idea.}
\begin{itemize}
\item Show that integration against bounded continuous test functions determines a probability measure on a Polish space (separating class).
\item Use regularity and a monotone class argument to pass from indicators of open/closed sets to bounded continuous functions, connecting set-based and function-based descriptions.
\item Conclude that the embedding $\mu\mapsto\big(\int f\,d\mu\big)_{f\in C_b(E)}$ is injective and metrizable via the Prokhorov metric framework.
\end{itemize}

\subsection*{Theorem 10.19 (Portmanteau theorem and Prokhorov metric)}
\textbf{Proof idea.}
\begin{itemize}
\item ($\Rightarrow$) If $\mu_n\Rightarrow\mu$, approximate indicators of open/closed sets by bounded continuous functions from below/above to get the liminf/limsup inequalities.
\item ($\Leftarrow$) Conversely, use the inequalities for all closed sets (or a convergence-determining class) to show convergence of integrals for bounded continuous functions via approximation.
\item Relate these conditions to the Prokhorov metric by showing $d_P(\mu_n,\mu)\to0$ iff the portmanteau inequalities hold.
\end{itemize}

\subsection*{Theorem 11.4 (Prokhorov’s theorem)}
\textbf{Proof idea.}
\begin{itemize}
\item Use tightness to extract a subsequence that is Cauchy in the Prokhorov metric by a diagonal argument over a countable determining class (e.g. bounded Lipschitz functions).
\item Show completeness/compactness: on a Polish space, tightness gives relative compactness of $\{\mu_n\}$ and any limit point is a probability measure.
\item Conversely, if $\{\mu_n\}$ is relatively compact, show it must be tight; otherwise mass would escape to infinity and prevent weakly convergent subsequences.
\end{itemize}

\subsection*{Theorem 11.5 (Prokhorov theorem on metric spaces)}
\textbf{Proof idea.}
\begin{itemize}
\item Adapt Prokhorov’s theorem from Polish spaces to general metric spaces by working on the completion or by restricting to the support of tight measures.
\item Show that tightness still prevents mass escape and yields sequential compactness in the weak topology (often for Radon measures).
\item Conversely, weak relative compactness forces tightness by testing on complements of large compact sets and using the portmanteau theorem.
\end{itemize}

\subsection*{Theorem 11.9 (Tightness criterion based on Arzel`a–Ascoli)}
\textbf{Proof idea.}
\begin{itemize}
\item Use Arzel\`a--Ascoli: relative compactness in $C([0,T];E)$ follows from (i) tightness at a fixed time (or at $0$) and (ii) uniform control of the modulus of continuity.
\item Translate (ii) into a probabilistic equicontinuity condition: for each $\varepsilon$ choose $\delta$ so that increments over intervals of length $\delta$ are small with high probability, uniformly in $n$.
\item Combine with a diagonal/compactness argument to get tightness of the laws on path space.
\end{itemize}

\subsection*{Theorem 11.16 (Kolmogorov–Chentsov)}
\textbf{Proof idea.}
\begin{itemize}
\item Assume moment bounds on increments: $\mathbb E[d(X_t,X_s)^\alpha]\le C|t-s|^{1+\beta}$.
\item Apply Markov’s inequality on dyadic grids and a Borel--Cantelli argument to show that, a.s., increments decay like $|t-s|^\gamma$ for some $\gamma<\beta/\alpha$.
\item Upgrade the dyadic control to all times by chaining (triangle inequality), yielding a modification with $\gamma$-H\"older continuous paths.
\end{itemize}

\subsection*{Theorem 11.17 (Kolmogorov criterion for tightness)}
\textbf{Proof idea.}
\begin{itemize}
\item Use Kolmogorov–Chentsov to get H\"older-continuous modifications under uniform increment moment bounds, which implies tightness in $C([0,T];E)$.
\item Alternatively, combine moment bounds with the Arzel\`a--Ascoli tightness criterion by controlling the modulus of continuity in probability.
\item Conclude that the family of process laws is tight on path space; hence weakly convergent subsequences exist.
\end{itemize}

\subsection*{Theorem 12.12 (Birkhoff (individual ergodic theorem)}
\textbf{Proof idea.}
\begin{itemize}
\item Define $\overline f=\limsup S_n$ and $\underline f=\liminf S_n$ and first show they are $T$-invariant using $S_n\circ T-S_n\to0$.
\item Apply Hopf’s maximal ergodic lemma to $h=f-g$ where $g$ is any invariant integrable function: it yields $\mathbb E[(f-g)\mathbf 1_{\{\sup_k (S_k(f)-g)>0\}}]\ge0$, giving the key comparison inequality.
\item Use this inequality twice to show $\int \overline f\,d\mathbb P\le \int g\,d\mathbb P$ for invariant majorants $g$ and $\int \underline f\,d\mathbb P\ge \int g\,d\mathbb P$ for invariant minorants; choosing $g=\underline f$ and $g=\overline f$ forces $\overline f=\underline f$ a.s.
\item Let $f^*=\lim S_n$; it is invariant, hence $\mathcal I$-measurable. Identify it by testing on invariant sets $A\in\mathcal I$: Hopf’s inequality applied to $f\mathbf 1_A$ gives $\int_A f^*=\int_A f$, which characterizes $\mathbb E[f\mid\mathcal I]$.
\item $L^1$ convergence follows by dominated convergence for bounded $f$ and truncation plus $L^1$-contraction for general $f\in L^1$.
\end{itemize}

\end{document}